\section{Формальная постановка и ядро Hybrid-INoT}
\label{sec:hybrid-core}

\subsection{Инженерная задача и наблюдаемый успех}
Инженерная задача задаётся тройкой $(x,\mathcal{C}_0,\mathcal{V})$: спецификацией $x$, фиксированным предоставленным контекстом $\mathcal{C}_0$ и внешней процедурой проверки $\mathcal{V}$. Результат включает программу-кандидат $y$ и запись её проверки. Генерация и проверка являются разными операциями: модель предлагает решение, а исходные исполняемые тесты определяют исход бенчмарка в контролируемом окружении. Текстовое утверждение о корректности не заменяет тесты. Для задач с репозиторием $\mathcal{C}_0$ обозначает раскрытый пакет найденных файлов, а не предположение о передаче всего репозитория.

Для протокола исполнения $a$ инженерная цель состоит в повышении вероятности успешного прохождения тестов при контроле ожидаемого числа токенов и денежной оценки генерации. Поэтому качество и ресурсы сообщаются совместно, без оптимизации непроверенного взвешенного показателя. Наблюдаемые программы, провалившие тесты, считаются неуспешными; недоступную генерацию или оценку нельзя превращать в успешный исход либо в вымышленный отказ тестов. Раздел~\ref{sec:heldout-methods} задаёт общий контроль оцениваемости и точный статистический исход.

\subsection{Разделение функций и внешняя проверка}
Hybrid-INoT задаёт в ответе модели три функции: планирование, реализацию и проверку требований; сама проверка программы выполняется после генерации. «Гибридность» означает сочетание инструкций в запросе с последующими исполняемыми тестами. Для этого не требуется обученный маршрутизатор. Исходная терминология planner--worker--critic соответствует исполняемым обозначениям \texttt{planner}, \texttt{implementer} и \texttt{reviewer}. Здесь изучается разделение предписанных функций, а не наличие независимо наблюдаемых скрытых агентов.

Исследуемое одновызовное ядро предписывает три операции: анализ требований и граничных случаев, реализацию по плану, затем проверку и исправление относительно требований. Один вызов модели получает все три инструкции и полный предоставленный контекст; он возвращает одну итоговую программу. Отдельные внутренние планы и оценки не наблюдаются. Поэтому концептуальное разделение функций не следует принимать за три измеренные подпрограммы. Точные запросы, включая требование к итоговому формату, приведены в приложении.

В основном эксперименте SR представляет эту ролевую одновызовную конфигурацию. SN является её абляцией обозначений ролей при одинаковых предложениях об операциях. MR и MN распределяют соответствующие операции по трём новым вызовам, а D полностью убирает поэтапную процедуру. Для всех пяти условий используется одинаковая последующая штатная оценка. Дизайн позволяет оценить, как заданные этапы соотносятся с результатом по сравнению с простой инструкцией и что меняется при распределении тех же операций по отдельным вызовам. У Hybrid-INoT нет преимущественного доступа к проверке.

\begin{figure}[htbp]\centering
\begin{tikzpicture}[
 x=1mm,y=1mm,
 box/.style={draw=black!80,align=center,font=\small,inner sep=2mm},
 flow/.style={-{Latex[length=2mm]},thick}]
\node[box,fill=gray!8,text width=100mm,minimum height=12mm] (input) at (0,0) {Задача и полный предоставленный контекст};
\node[box,text width=43mm,minimum height=54mm] (single) at (-28,-40) {\textbf{SR: один вызов}\\[3mm]В одном запросе:\\[1mm]планирование;\\реализация;\\проверка требований\\[3mm]Внутренние этапы\\не измеряются};
\node[box,text width=43mm,minimum height=15mm] (m1) at (28,-20) {\textbf{MR: вызов 1}\\Составить план};
\node[box,text width=43mm,minimum height=15mm] (m2) at (28,-40) {\textbf{MR: вызов 2}\\Написать программу\\{\footnotesize Все предыдущие ответы}};
\node[box,text width=43mm,minimum height=15mm] (m3) at (28,-60) {\textbf{MR: вызов 3}\\Проверить и исправить\\{\footnotesize Все предыдущие ответы}};
\node[box,text width=100mm,minimum height=14mm] (candidate) at (0,-84) {\textbf{Извлечение из последнего ответа}\\Одна программа; неверный формат даёт пустого кандидата};
\node[box,fill=gray!8,text width=100mm,minimum height=18mm] (test) at (0,-110) {\textbf{Внешняя проверка после генерации}\\Исходные тесты бенчмарка и общие контрольные проверки};
\node[box,text width=100mm,minimum height=20mm] (outcome) at (0,-139) {\textbf{Итог проверки}\\Успех, неуспех или недоступная оценка\\Расход токенов учитывается отдельно};
\draw[flow,dashed] (input.south) -- ++(0,-3) -| (single.north);
\draw[dashed,thick] (input.east) -- (61,0) -- (61,-60);
\foreach \call in {m1,m2,m3} {\draw[flow,dashed] (61,0 |- \call.east) -- (\call.east);}
\node[rotate=90,font=\footnotesize] at (65,-37) {Та же задача и контекст};
\draw[flow] (m1) -- (m2);
\draw[flow] (m2) -- (m3);
\draw[flow] (single.south) -- ++(0,-5) -| ([xshift=-28mm]candidate.north);
\draw[flow] (m3.south) -- ++(0,-5) -| ([xshift=28mm]candidate.north);
\draw[flow] (candidate) -- (test);
\draw[flow] (test) -- (outcome);
\end{tikzpicture}

\caption{Два варианта генерации: SR задаёт три функции одним запросом, MR распределяет их по трём вызовам. Стрелки между вызовами MR означают передачу полных предыдущих ответов; задача и контекст доступны в каждом вызове. Показан путь полученного ответа. Незавершённая генерация учитывается отдельно по алгоритму~\ref{alg:hybrid-core}. Все пять условий из табл.~\ref{tab:design} используют одну процедуру внешней оценки, без передачи результатов тестов обратно модели.}
\label{fig:hybrid-core}
\end{figure}
\FloatBarrier

\subsection{Состояние, исполнение и остановка}
Наблюдаемое состояние исполнения задаётся как $\sigma_j=(x,\mathcal{C}_0,a,j,\mathcal{H}_j)$, где $a$ фиксирует назначенное условие, $j$ --- номер вызова, а $\mathcal{H}_j$ содержит все полные предыдущие выведенные ответы. Конфигурация модели внутри сравнения фиксирована. Для D, SR и SN число вызовов $n_a=1$; для MR и MN $n_a=3$. В многовызовном условии $\mathcal{H}_j$ передаётся целиком. Предоставленный контекст не изменяется: $\widehat{\mathcal{C}}=\mathcal{C}_0$. Условие назначается до генерации и потому не является политикой маршрутизации по результату.

\begin{algorithm}[htbp]
\caption{Генерация программы и проверка результата}
\label{alg:hybrid-core}
\small
\textbf{Вход:} задача, полный предоставленный контекст и заранее назначенное условие.\par
\textbf{Выход:} полученная программа, исход проверки и известный расход ресурсов.\par\medskip
\textbf{I. Генерация}
\begin{enumerate}\setlength{\itemsep}{2pt}\setlength{\parsep}{0pt}
\item Определить число вызовов: один для D, SN или SR; три для MN или MR.
\item Выполнить вызовы последовательно. В каждый передать задачу, неизменный контекст и заданную инструкцию. Во второй и третий вызовы добавить все полные предыдущие ответы.
\item Для каждого вызова сохранить запрос, полный ответ, конечный статус и счётчики токенов.
\item Если вызов не завершён или учёт ресурсов не подтверждён, сохранить доступные сведения и остановить это назначение как недоступное. Не повторять его автоматически.
\item Из последнего завершённого ответа извлечь одну программу. Если формат неверен, сохранить пустого кандидата.
\end{enumerate}
\textbf{II. Проверка после окончания генерации}
\begin{enumerate}\setcounter{enumi}{5}\setlength{\itemsep}{2pt}\setlength{\parsep}{0pt}
\item Проверить каждого полученного кандидата исходными тестами бенчмарка.
\item Сопоставить проверенную программу, отчёт оценщика и результаты общих контрольных проверок.
\item Записать успех, неуспех или недоступный исход. Сохранить все известные показатели ресурсов; не возвращать результаты тестов модели.
\end{enumerate}
\end{algorithm}

Алгоритм~\ref{alg:hybrid-core} описывает этапы генерации и оценки эксперимента; штатный оценщик запускается после генерации и не является инструментом внутри вызова модели. Лимит времени и конечное $n_a$ ограничивают попытки генерации. Это правило остановки, а не теорема о сходимости или корректности. Отклонённое завершение не вызывает автоматическую повторную попытку той же ячейки. Раскрытое продолжение касается только назначений, в которых ни один вызов не начинался.

Конфигурация фиксирует дополнительные механизмы исходной архитектуры: отбор контекста тождественный, режим назначен, итоговый кандидат один, повторные генерации по проверке отключены. Так вопрос о заданных функциях можно исследовать, не смешивая его с компрессией, малой проверяющей моделью и адаптивными повторами. Этим же определяется объём вывода: эксперимент оценивает данное ядро, а не все возможные варианты применения Hybrid-INoT.
